\begin{pastexam}[label={px:ntut-108-ci-06}]{北科大資工所 108 學年度 第 4 題}
（是非）

(6 pts) Let $A_1,\ldots,A_5$ be matrices with dimension $2\times1$, $1\times3$, $3\times2$, $2\times4$,
$4\times1$, respectively. Let $C_{i,j}$ be the smallest number of scalar multiplications needed for
computing the matrix product $A_iA_{i+1}\ldots A_j$, assuming that multiplying an $r\times s$ matrix
and an $s\times t$ matrix takes $rst$ scalar multiplications. Mark by TRUE or FALSE each of the
following statements:
\begin{enumerate}[label=(\arabic*)]
\item (2 pts) $C_{3,4}$ is computed before $C_{2,5}$.
\item (2 pts) $C_{1,4}=22$.
\item (2 pts) $C_{1,5}$ is derived from $C_{1,2}$ and $C_{3,5}$.
\end{enumerate}
\end{pastexam}

\begin{pastsoln}
答案：(1) TRUE；(2) TRUE；(3) FALSE。

$A_i$ 是 $p_{i-1} \times p_i$，維度序列 $\langle p_0, \ldots, p_5 \rangle = \langle 2, 1, 3, 2, 4, 1 \rangle$。
照資結書第~18.2~節的遞迴式
\[
C_{i,i} = 0, \qquad
C_{i,j} = \min_{i \le k < j} \bigl( C_{i,k} + C_{k+1,j} + p_{i-1} p_k p_j \bigr) \quad (i < j),
\]
依鏈長由短到長填表。 每一格列出所有 $k$：

\begin{center}
\small
\begin{tabular}{lll}
\toprule
格子 & 各個 $k$ 的值 & 最小值 \\
\midrule
$C_{1,2}$ & $k{=}1$：$0 + 0 + 2 \cdot 1 \cdot 3 = 6$ & $6$ \\
$C_{2,3}$ & $k{=}2$：$0 + 0 + 1 \cdot 3 \cdot 2 = 6$ & $6$ \\
$C_{3,4}$ & $k{=}3$：$0 + 0 + 3 \cdot 2 \cdot 4 = 24$ & $24$ \\
$C_{4,5}$ & $k{=}4$：$0 + 0 + 2 \cdot 4 \cdot 1 = 8$ & $8$ \\
\midrule
$C_{1,3}$ & $k{=}1$：$0 + 6 + 2 \cdot 1 \cdot 2 = 10$；$k{=}2$：$6 + 0 + 2 \cdot 3 \cdot 2 = 18$ & $10$（$k{=}1$） \\
$C_{2,4}$ & $k{=}2$：$0 + 24 + 1 \cdot 3 \cdot 4 = 36$；$k{=}3$：$6 + 0 + 1 \cdot 2 \cdot 4 = 14$ & $14$（$k{=}3$） \\
$C_{3,5}$ & $k{=}3$：$0 + 8 + 3 \cdot 2 \cdot 1 = 14$；$k{=}4$：$24 + 0 + 3 \cdot 4 \cdot 1 = 36$ & $14$（$k{=}3$） \\
\midrule
$C_{1,4}$ & $k{=}1$：$0 + 14 + 8 = 22$；$k{=}2$：$6 + 24 + 24 = 54$；$k{=}3$：$10 + 0 + 16 = 26$ & $22$（$k{=}1$） \\
$C_{2,5}$ & $k{=}2$：$0 + 14 + 3 = 17$；$k{=}3$：$6 + 8 + 2 = 16$；$k{=}4$：$14 + 0 + 4 = 18$ & $16$（$k{=}3$） \\
\midrule
$C_{1,5}$ & $k{=}1$：$0 + 16 + 2 = 18$；$k{=}2$：$6 + 14 + 6 = 26$； & \\
 & $k{=}3$：$10 + 8 + 4 = 22$；$k{=}4$：$22 + 0 + 8 = 30$ & $18$（$k{=}1$） \\
\bottomrule
\end{tabular}
\end{center}
後三段每個 $k$ 的第三項分別是 $p_{i-1} p_k p_j$，例如 $C_{1,4}$ 的 $k = 1$ 是 $2 \cdot 1 \cdot 4 = 8$。

\textbf{(1) TRUE。} 由下而上填表時依鏈長由短到長（資結書第~18.2~節），$C_{3,4}$ 鏈長 $2$，
$C_{2,5}$ 鏈長 $4$，前者先算。 不只是這個順序：$C_{2,5}$ 的 $k = 2$ 那一項要用 $C_{3,5}$，而
$C_{3,5}$ 的 $k = 4$ 那一項要用 $C_{3,4}$，所以任何正確的計算順序（包括由上而下的記憶化，
看的是值算完的先後）都必須先算好 $C_{3,4}$，才能算完 $C_{2,5}$。

\textbf{(2) TRUE。} 見表：$C_{1,4} = \min(22, 54, 26) = 22$，在 $k = 1$ 取到，也就是
$A_1 (A_2 A_3 A_4)$，而 $A_2 A_3 A_4$ 用 $(A_2 A_3) A_4$ 花 $14$ 次。

\textbf{(3) FALSE。} 「由 $C_{1,2}$ 和 $C_{3,5}$ 得到」是在 $k = 2$ 切開，那一項是
$6 + 14 + 2 \cdot 3 \cdot 1 = 26$。 但最小值是 $k = 1$ 的 $18$：$C_{1,1} + C_{2,5} + 2 \cdot 1 \cdot 1$，
也就是 $A_1 (A_2 A_3 A_4 A_5)$。 所以 $C_{1,5} = 18$ 是由 $C_{1,1}$ 和 $C_{2,5}$ 得到的。
若把「derived from」讀成「計算時有用到」，$k = 2$ 那一項確實有讀 $C_{1,2}$、$C_{3,5}$，但
$k = 1, 3, 4$ 的項也都讀了，那樣讀這句話就沒有意義；照一般的意思（最佳解從哪兩個子問題組成），
答 FALSE。
\end{pastsoln}
